\caption{\textbf{Testing structured vs.\ random mixed selectivity beyond second-order statistics.}
Feature vectors $X$ (54{,}719 neurons $\times$ 1872 bins; concatenated z-scored PETHs of 21 trial-type/alignment windows).
\textbf{a,} Correlation matrix of the coefficients $C = XV^\top$ on the 100 k-means cluster means $V$ (as in Fig.~4); rows/columns in hierarchical order. Its off-diagonal pattern equals that of white noise projected onto $V$ ($r = 0.91$), and PC1 holds 71\% of the coefficient variance: the pattern is set by the non-orthogonal basis.
\textbf{b,} Same for an orthonormal basis $W = (VV^\top)^{-1/2}V$ (L\"owdin). White noise now gives the identity ($r = 0.00$), and the remaining structure (mean $|r| = 0.17$, PC1 24\%) is the covariance of the responses.
\textbf{c,} Mean $|$off-diagonal correlation$|$ for the real coefficients, white noise, the axis-marginal null of Fig.~4 ($\approx 0$ by construction) and a covariance-matched Gaussian (identical to the real data by construction). Neither null tests anything beyond second order.
\textbf{d,} Cumulative variance of the PCs of $X$ (10 PCs: 40\%; 20 PCs: 47\%).
\textbf{e,} ePAIRS-style test in the top 10 PCs: distribution of each neuron's mean angle to its 5 nearest neighbours (neighbours from other probe insertions only), real (black) vs a covariance-matched Gaussian (blue), a PC-marginal null (each PC score permuted independently; orange) and a positive control, a 25-component Gaussian mixture fitted by k-means to the same data (pink, dashed). AUC $=P(\theta_\mathrm{real}<\theta_\mathrm{null})$; $\Delta$ is the difference in median angle.
\textbf{f,} Same for the 100 orthonormalised coefficients $XW^\top$, with the axis-marginal null of Fig.~4 (yellow), which separates trivially (AUC 0.98) because it also destroys the covariance.
\textbf{g,} Gap statistic of k-means ($\log W_k$ of the null minus the real one; mean of 10 null draws, s.d.\ $\le 0.003$) in the top 10 (solid) and 20 (dotted) PCs, vs the Gaussian, PC-marginal and elliptical (Gaussian directions, real Mahalanobis radii) nulls; pink, the mixture control vs the Gaussian null.
\textbf{h,} Mean silhouette (10{,}000 neurons) vs $k$: real data lie slightly above all nulls but decrease monotonically, with no preferred number of clusters (the mixture control shows a bump near $k=20$).
\textbf{i,} Correlations across neurons between 11 task-selectivity indices (window-mean contrasts of condition-averaged responses; lower triangle). Upper triangle: correlation predicted by trial-averaging noise shared between contrasts that use the same windows (estimated from lag-14-ms differences of the PETHs). Side (stimulus = choice on correct trials; not dissociable) and block indices correlate negatively ($-0.31$) because of unequal trial numbers in prior-congruent and -incongruent windows (predicted $-0.20$); correlations between epochs (e.g.\ block pre-stim vs at stim, 0.63; side at stim vs before movement, 0.74) reflect persistence of the same variable.
\textbf{j,} ePAIRS on the 11-dimensional selectivity vectors, all neurons (solid) and the 20\% with the largest selectivity norm (dashed; same selection applied to each null).
\textbf{k,} Effect sizes (AUC, 5 nearest neighbours, cross-probe) for each dataset and null; 0.5 = no excess angular clustering. Real responses are more clustered than any second-order-matched null (AUC 0.58--0.66), more so than the 25-cluster mixture fitted to them (0.58), but without discrete clusters (g, h); task-selectivity vectors are indistinguishable from a Gaussian (AUC 0.50--0.53), and the excess in the most selective neurons (0.68 vs Gaussian) shrinks to 0.57 when the PC marginals are preserved.}
\label{fig:selectivity_test}
